\scanpage{16}
Polygône régulier (ou réalisation régulière d'un polygône combinatoire) : définitions possibles (\unsure{cotés|côtés} égaux, angles égaux, \inserted{\deleted{\unsure{il suffit|il faut|il y avait}} les deux,} transitivité groupe des symétries ...) Existence, unicité modulo homothétie : tout isom. combinatoire est induit par une unique similitude.

\figure{\begin{tikzpicture}[x=.8cm,y=.8cm]
\draw (0,0) -- (1.14,.66);
\draw (1.18,.70) -- (1.55,1.97);
\end{tikzpicture}}

Notion d'enveloppe affine (ou vectorielle) d'un polygône combinatoire ... \unsure{sens} précisé

\noindent\hfil\rule{.24\linewidth}{.35pt}\hfil\par
Fonctorialité de $X\times_G Y$ \unsure{p.r.} à $(\unsure{G|X},X,Y)$ ...
\par\noindent\hfil\rule{.27\linewidth}{.35pt}\hfil\par

\[
\begin{aligned}
\unsure{P|\overline{P}|T}\times_G E={}^{P}\!\deleted{\unsure{B}}E=F
&\qquad\qquad &P&\quad G\text{-torseur}\\
&&E&\quad G\text{-ensemble}
\end{aligned}
\]

$P\longrightarrow\operatorname{Bij}(E,F)$ compatible aux opérations de $G$ à droite

\uline{NB} Si $\unsure{G}$ opère fidèlement sur $E$, $G$ opère ``sans pts fixes'' sur $\operatorname{Bij}(E,F)$ --- i.e. $\deleted{\unsure{y}}\unsure{x|y}.g=\unsure{x|u}$ $(g\in G,\ x\in B)$ implique $g=1$, ou encore les $G_x$ $(x\in B)$ sont réduits à $\{1\}$, ou encore $B$ est une somme de torseurs ... Dans ce cas\inserted{ et dans ce cas seulement}, $P\mathbin{\text{---}}\operatorname{Bij}(E,F)$ est injectif

\[
P\hookrightarrow\operatorname{Bij}(E,F)
\]

% slanted
\margin{\tikz[baseline=(side.base)]{
\node[inner sep=1pt] (side) {\shortstack[l]{$(G,E)$-structure\\sur ens $F$}};
\draw (side.south west) .. controls +(1em,-.3em) and +(-.35em,.7em) .. ([xshift=-.04em,yshift=-.5em]side.south)
 .. controls +(.3em,.8em) and +(-.8em,-.3em) .. (side.south east);
\node[anchor=north east,inner sep=1pt] at ([xshift=-.2em,yshift=-.25em]side.south) {\deleted{\illegible{1}}};
\draw ([yshift=-.57em]side.south) -- ++(0,-2em) .. controls +(0,-.4em) and +(-.3em,0) .. ++(.55em,-.5em);}}
\noindent$\left[\begin{minipage}{.84\linewidth}
En tous les cas, $F\inserted{={}^{P}\!E}$ est muni de la structure suivante : donnée d'une orbite de $G$ dans $\operatorname{Bij}(E,F)$.
\end{minipage}\right.$

{\footnotesize [On verra que (si $G$ opère $\unsure{\text{fid}^{\mathrm{t}}}$ sur $E$) la donnée d'un \deleted{\unsure{ens|d'un}} \deleted{\unsure{$(G,E)$-st}} $\unsure{F}$ avec $(G,E)$-structure \unsure{équivaut à|équiv. à} celle d'un \unsure{$G$-torseur} \unsure{et applic$^{\mathrm{tn}}$}.]}

\uline{Corollaire} \unsure{de|du} $\unsure{\big[\,\mathbin{\text{---}}|\big[\,\longrightarrow}$ \unsure{toute structure}
% archive sheet number: 15

